\scanpage{124}
\mspage{61}
considérons les images inverses de $T_d$ et $T_{d'}$ sur $D_{d,d'}$, qu'on va noter
\deleted{${\unsure{T|\mathbb Z}_d}_{D_{d,d'}}$}
\[
(193)\quad\left\{\begin{aligned}
T(d,d')_d&=(T_d)_{D_{d,d'}}&&\begin{gathered}\text{syst. local}\\\text{de rang }d\text{ sur }D_{d,d'}\end{gathered}\\
T(d^{\unsure{*}},\unsure{d'|d})_{d'}&=(T_{d'})_{D_{d^{\unsure{*}},d'}}&&\begin{gathered}\text{syst. local de}\\\text{rang }d'\text{ sur }D_{d,d'}\end{gathered}
\end{aligned}\right.
\]
on a alors un hom canonique surjectif (dit ``tautologique'')
\[
(194)\qquad T(d,d')_d\mathrel{\tikz[baseline=-.5ex,x=1em,y=1ex]{\draw (0,0)--(1.81,0);\draw (1.8,.15)--(1.82,-.15);}}T(d.d')_{d'}
\]
qui, par image inverse sur $D^!_{d,d'}=D^!_d\mathbin\times_{D_d}D_{d'}$ (cf diagramme can.
\diagram{\begin{tikzcd}[ampersand replacement=\&,column sep=large,row sep=large]
\&\&\mathcal P_{d'}\raisebox{.6ex}{\((I_d)\)}{}_{D^!_d}\simeq D^!_{d,d'}\arrow[dl]\arrow[dr,dash,"\mathfrak S_{d^{\unsure{*}}}",start anchor=south east,to path={-- (\tikztotarget.north west)\tikztonodes -- ([xshift=-1pt,yshift=2.3pt]\tikztotarget.north west)}]\arrow[rr,dashed]\&\&D^!_{d'}\arrow[dd,"\mathfrak S_{d'}"]\\
(195)\&D^!_d\arrow[dr,dash,"\mathfrak S_d"]\&\&D_{d,d'}\arrow[dl,dash,start anchor=south west,to path={-- (\tikztotarget.north east)\tikztonodes -- ([xshift=-1pt,yshift=2.3pt]\tikztotarget.north east)}]\arrow[dr]\&\\
\&\&D_d\&\&D_{d'}
\end{tikzcd}}
s'identifie à l'hom de projection
\[
(196)\qquad\mathbb Z^d\mathrel{\tikz[baseline=-.5ex,x=1em,y=1ex]{\draw (0,0)--(1.81,0);\draw (1.8,.15)--(1.82,-.15);}}\mathbb Z^{d'}
\]
de $\mathbb Z^d$ sur un produit partiel défini par une partie de cardinal $d'$ de l'ens $I_d$ en exposant : $\mathbb Z^d\simeq\mathbb Z^{I_d}$\unsure{.|,}
\margin{123}
